\subsection{Критерии Сильвестра}
\begin{theorem}[Критерий Сильвестра положительно определенной формы]
    Квадратичная форма положительно определена тогда и только тогда, когда все главные миноры её матрицы положительны.
\end{theorem}
\begin{proof}
Пусть задана квадратичная форма $k(x)$, её порождающая билинейная форма $b(x,y)$ с матрицей $B$.
\begin{enumerate}
    \item[$\Rightarrow$] Если $k(x)$ положительно определена, то в \textit{любом базисе} диагональные элементы матрицы $B$ положительны ($b_{ii}=b(e_i,e_i)=k(e_i)>0$). Приведем $B$ элементарными преобразованиями к диагональному виду. При этом мы не встретим "особого случая" из теоремы 9.1, так как диагональные элементы положительны. Значит в ходе приведения выполняются следующие условия, следующие из алгоритма теоремы 9.1 в "основном случае":
    \begin{itemize}
        \item Для строк: вычитаются только строки, находящиеся выше.
        \item Для столбцов: вычитаются только столбцы, находящиеся левее.
    \end{itemize}
    \[
\begin{tikzpicture}[baseline=(current bounding box.center)]
    % Первая матрица с круглыми скобками
    \matrix (m1) [matrix of math nodes, nodes in empty cells,
        left delimiter={(}, right delimiter={)}] {
        b_{11} & b_{12} & \cdots & b_{1n} \\
        b_{21} & b_{22} & \cdots & b_{2n} \\
        \vdots & \vdots & \ddots & \vdots \\
        b_{n1} & b_{n2} & \cdots & b_{nn} \\
    };
    
    % Линии для первой матрицы
    \draw[gray, thick] (m1-1-1.north west) rectangle (m1-1-1.south east);
    \draw[gray, densely dashed] (m1-1-1.north west) rectangle (m1-2-2.south east);
    \draw[gray, loosely dashed] (m1-1-1.north west) rectangle (m1-3-3.south east);

    % Вторая матрица с круглыми скобками
    \matrix (m2) [matrix of math nodes, nodes in empty cells,
        left delimiter={(}, right delimiter={)},
        right=3cm of m1.north east, anchor=north west] {
        b_{11} & 0 & \cdots & 0 \\
        0 & b_{22}' & \cdots & b_{2n}' \\
        \vdots & \vdots & \ddots & \vdots \\
        0 & b_{n2}' & \cdots & b_{nn} \\
    };
    
    % Линии для второй матрицы
    \draw[gray, thick] (m2-1-1.north west) rectangle (m2-1-1.south east);
    \draw[gray, densely dashed] (m2-1-1.north west) rectangle (m2-2-2.south east);
    \draw[gray, loosely dashed] (m2-1-1.north west) rectangle (m2-3-3.south east);

    % Стрелка между матрицами (с отступом)
    \draw[-Latex, thick, shorten >=20pt, shorten <=20pt] 
        (m1.east) -- node[above, pos=0.5] {$S^TBS$} (m2.west);
\end{tikzpicture}
\]
    А из этих условий следует следующий факт: данные элементарные преобразования матрицы $B$ являются элементарными для всех её подматриц, соответствующих главным минорам, а значит главные миноры не меняются при смене базиса.
    
    После приведения к диагональному виду получается матрица
    \[
    B'=diag(b_{11}, b_{22}', \dots, b_{nn}'), \quad b_{ii}'>0
    \]
    А все главные миноры $\Delta_i$ являются определителями диагональных матриц с положительными элементами, следовательно, они положительны
    \[
    \Delta_1=b_{11}, \quad \Delta_2=b_{11}b_{22}', \quad \dots, \quad \Delta_n=b_{11}b_{22}'\dots b_{nn}'
    \]
    \item[$\Leftarrow$] Если все главные миноры $\Delta_i$ положительны. Так как $\Delta_1 = b_{11}>0$, то мы можем сделать хотя бы одну итерацию "основного случая" алгоритма приведения матрицы $B$ к диагональному виду (аналогично п.1 при этом сохраняются главные миноры). Тогда получим 
    \[
\begin{tikzpicture}[baseline=(current bounding box.center)]
    % Первая матрица с круглыми скобками
    \matrix (m1) [matrix of math nodes, nodes in empty cells,
        left delimiter={(}, right delimiter={)}] {
        b_{11} & b_{12} & \cdots & b_{1n} \\
        b_{21} & b_{22} & \cdots & b_{2n} \\
        \vdots & \vdots & \ddots & \vdots \\
        b_{n1} & b_{n2} & \cdots & b_{nn} \\
    };
    
    % Линии для первой матрицы
    \draw[gray, thick] (m1-1-1.north west) rectangle (m1-1-1.south east);
    \draw[gray, densely dashed] (m1-1-1.north west) rectangle (m1-2-2.south east);
    \draw[gray, loosely dashed] (m1-1-1.north west) rectangle (m1-3-3.south east);

    % Вторая матрица с круглыми скобками
    \matrix (m2) [matrix of math nodes, nodes in empty cells,
        left delimiter={(}, right delimiter={)},
        right=3cm of m1.north east, anchor=north west] {
        b_{11} & 0 & \cdots & 0 \\
        0 & b_{22}' & \cdots & b_{2n}' \\
        \vdots & \vdots & \ddots & \vdots \\
        0 & b_{n2}' & \cdots & b_{nn} \\
    };
    
    % Линии для второй матрицы
    \draw[gray, thick] (m2-1-1.north west) rectangle (m2-1-1.south east);
    \draw[gray, densely dashed] (m2-1-1.north west) rectangle (m2-2-2.south east);
    \draw[gray, loosely dashed] (m2-1-1.north west) rectangle (m2-3-3.south east);

    % Стрелка между матрицами (с отступом)
    \draw[-Latex, thick, shorten >=20pt, shorten <=20pt] 
        (m1.east) -- node[above, pos=0.5] {$S^TBS$} (m2.west);
\end{tikzpicture}
\]
Но $\Delta_2=b_{11}b_{22}'>0$, а значит $b_{22}'>0$ и мы можем сделать следующую итерацию. После $k$-ой итерации получаем $\Delta_k=\Delta_{k-1} b_{kk}'>0$ и $\Delta_{k-1}>0$, откуда $b_{kk}'>0$ и следующая итерация возможна. После $n$-ого шага получаем диагональную матрицу с положительными элементами
\[
B'=diag(b_{11}, b_{22}', \dots, b_{nn}'), \quad b_{ii}'>0
\]
Откуда следует положительная определенность квадратичной формы.
\end{enumerate}
\end{proof}
\textit{Замечание}: главные миноры (определители выделенных пунктиром в подматриц) иногда называют угловыми минорами.

\textit{Замечание}: в доказательстве использовался факт о том, что из $k(x)>0$ следует, что $b_{ii}>0$. Обратное неверно, контрпример: $k(x_1,x_2)=x_1^2-2x_1x_2+x_2^2$, $b_{11}=b_{22}>0$ и $k(1,1)=0$.

\begin{theorem}[Критерий Сильвестра отрицательно определенной формы]
    Квадратичная форма отрицательно определена тогда и только тогда, когда знаки главных миноров чередуются начиная с "$-$".
\end{theorem}
\begin{proof}
    Форма $k(x)$ отрицательно определена $\Leftrightarrow$ форма $(-1)k(x)$ положительно определена. Значит по критерию для положительно определенной формы все главные миноры $(-1)^m\Delta_m>0$. Откуда получаем, что главные миноры $\Delta_m$ формы $k(x)$ имеют чередующиеся знаки, причем $(-1)\Delta_1>0$, а значит $\Delta_1 <0$ и чередование начинается с "$-$".    
\end{proof}
\textit{Замечание}: из неотрицательности главных миноров не следует полуопределенность квадратичной формы. Контрпример: форма с матрицей $diag(0, -1, 1)$.

Пусть $k(x)$ -- квадратичная форма, определенная на $L$, где $\dim L =n\geqslant 2$. $A$ -- матрица формы $k(x)$.
\begin{lemma}[для матанализа]
    Если $\Delta_2=\det \begin{pmatrix}
        a_{11} & a_{12} \\
        a_{21} & a_{22}
    \end{pmatrix} < 0$, то $k(x)$ знаконеопределена.
\end{lemma}
\begin{proof}
    Покажем, что $k(x)$ знаконеопределена на $L^*=\langle e_1, e_2\rangle$. На пространстве $L^*$ у $k(x)$ матрица имеет вид $B=\begin{pmatrix}
        a_{11} & a_{12} \\
        a_{21} & a_{22}
    \end{pmatrix}$. Приведем ее к каноническому виду с помощью элементарного преобразования: $B'=S^TBS$. По условию $\det B < 0$, следовательно $\det B (\det S)^2 <0 \Leftrightarrow \det B'<0$. Отсюда $B'$ имеет вид $\begin{pmatrix}
        1 & 0 \\
        0 & -1
    \end{pmatrix}$ или $\begin{pmatrix}
        -1 & 0\\
        0 & 1
    \end{pmatrix}$. Значит $k(x)$ знаконеопределена.
\end{proof}

\section{Евклидово пространство}
% \subsection{Евклидово пространство}
\begin{definition}
    Вещественное линейное пространство $\mathcal{E}$ называется евклидовым, если на нем введена операция скалярного умножения [$\forall x,y \in \mathcal{E} \hookrightarrow (x,y)\in \R$], удовлетворяющая следующим свойствам
    \begin{enumerate}
        \item $\forall x,y \in \mathcal{E} \hookrightarrow(x,y)=(y,x)$
        \item $\forall x, y, z \in \mathcal{E} \ \forall\alpha, \beta \in \R \hookrightarrow (\alpha x+ \beta y, z)=\alpha(x,z)+\beta(y,z)$
        \item $\forall x\neq 0 \in \mathcal{E} \hookrightarrow (x,x) > 0$
    \end{enumerate}
\end{definition}
Менее формальная формулировка: вещественное линейное пространство $\mathcal{E}$ евклидово, если на нем задана положительно определенная квадратичная форма.

\begin{definition}
    Длина вектора $x\in \mathcal{E}$ задается формулой $$|x|=\sqrt{(x,x)}$$
\end{definition}
\begin{definition}
    Угол между векторами $x,y\neq 0 \in \mathcal{E}$ задается формулой 
    \[
    \cos \Phi = \frac{(x,y)}{|x||y|}, \quad \Phi \in [0, \pi]
    \]
\end{definition}
Корректность данного определения угла следует из неравенства Коши-Буняковского-Шварца.
\begin{theorem}[Неравенство КБШ]
    Для любых векторов $x,y\in \mathcal{E}$ справедливо неравенство
    \[
        (x,y)^2\leqslant(x,x)(y,y)
    \]
\end{theorem}
Доказательство данной теоремы будет приведено \hyperref[KBH proof]{далее}.
\begin{theorem}[Неравенство треугольника]
    Для любых векторов $x,y\in \mathcal{E}$ справедливо неравенство
    \[
    |x+y|\leqslant|x|+|y|
    \]
\end{theorem}
\begin{proof}
    В силу линейности и симметричности скалярного произведения
    \[
    |x+y|^2=(x+y, x+y)=(x,x)+2(x,y)+(y,y)
    \]
    По неравенству КБШ
    \[
    |x+y|^2=(x,x)+2(x,y)+(y,y) \leqslant (x,x)+2\sqrt{(x,x)(y,y)}+(y,y)=|x|^2+2|x||y|+|y|^2=(|x|+|y|)^2
    \]
    Откуда получаем $|x+y|\leqslant|x|+|y|$. Причем равенство достигается тогда же, когда и в КБШ.
\end{proof}
\begin{definition}
    Векторы $x,y \in \mathcal{E}$ называются ортогональными, если $(x,y)=0$.
\end{definition}
Из определения следует, что $x,y$ -- ортогональны $\Rightarrow\left[
\begin{aligned}
x &= 0 \\
y &= 0\\
\text{угол}&=\pi/2
\end{aligned}
\right.$.
\begin{proposition}
    Только $\overline{0}$ вектор ортогонален любому вектору из $\mathcal{E}$.
\end{proposition}
\begin{proof}
    Если вектор $x$ ортогонален любому вектору, то он ортогонален сам себе
    \[
    (x, x)=0 \Rightarrow x=0
    \]
\end{proof}

\subsection{Матрица Грама}
Для скалярного произведения в базисе  $e^{(n)}$ так же верна формула
\[
(x,y)=(x_1, \dots, x_n)\begin{pmatrix}
    (e_1,e_1) & \dots & (e_1,e_n) \\
    \vdots & \ddots & \vdots \\
    (e_n,e_1) & \dots & (e_n, e_n)
\end{pmatrix}\begin{pmatrix}
    y_1 \\
    \vdots \\
    y_n
\end{pmatrix}=(x_1, \dots, x_n)(\textbf{Г})\begin{pmatrix}
    y_1 \\
    \vdots \\
    y_n
\end{pmatrix}
\]
где матрица \textbf{Г} называется матрицей Грама базиса $e^{(n)}$.

\begin{theorem}
    Пусть задан произвольный набор векторов $f_1, \dots, f_n$. Тогда $\det \textbf{Г}_f\geqslant 0$, причем
    \begin{itemize}
        \item $\det \textbf{Г}_f > 0$, если набор ЛНЗ.
        \item $\det \textbf{Г}_f =0$, если набор ЛЗ.
    \end{itemize}
\end{theorem}
\begin{proof}
    Если набор ЛНЗ, то он образует базис. Тогда по критерию Сильвестра для положительной формы (квадратичная форма, порожденная скалярным произведением положительно определена) $\det \textbf{Г}_f >0$.

    Если набор ЛЗ, то существует нетривиальная линейная комбинация $\alpha_1, \dots, \alpha_n$ такая, что $\alpha_1f_1+\dots + \alpha_nf_n=0$. Последовательно скалярно умножим справа это равенство на $f_i, i\in\overline{1,n}$, получим систему
    \[
    \left\{
    \begin{aligned}
    &\alpha_1(f_1, f_1)+\dots+\alpha_n(f_n,f_1)=0\\
    &\alpha_1(f_1, f_2)+\dots+\alpha_n(f_n,f_2)=0\\
    &\dots\\
    &\alpha_1(f_1, f_n)+\dots+\alpha_n(f_n,f_n)=0
    \end{aligned}
    \right.
    \]
    Это однородная система с матрицей \textbf{Г} у которой есть нетривиальное решение $\alpha^{(n)}$. Следовательно, $\det \textbf{Г}_f=0$ (по теореме Крамера).
\end{proof}
\textit{Следствие}:\label{KBH proof} Докажем неравенство КБШ: применим данную теорему к набору $x, y$
\[
\det \begin{pmatrix}
    (x,x) & (x,y)\\
    (y,x) & (y,y)
\end{pmatrix} \geqslant 0 \Rightarrow (x,y)^2\leqslant(x,x)(y,y)
\]
Причем равенство достигается тогда и только тогда, когда $x,y$ -- ЛЗ.

\textit{Замечание}: при замене базиса $\textbf{Г}'=S^T\textbf{Г}S$.
\subsection{Ортогональный базис}
\begin{definition}
    Базис является ортогональным, если 
    \[
    (e_i, e_j)=0, \quad i \neq j
    \]
\end{definition}
\begin{definition}
    Базис является ортонормированным, если 
    \[
    (e_i, e_j)=\left\{
    \begin{aligned}
    0, \quad i\neq j\\
    1, \quad i = j
    \end{aligned}
    \right. \quad \forall i, j \in \overline{1,n}
    \]
\end{definition}

В ОНБ матрица Грама вырождается в единичную: $\textbf{Г}=E$. Соответственно формула скалярного произведения принимает простой вид
\[
(x,y)
 = \sum_{i=1}^nx_iy_i
\]
\begin{proposition}
    Пусть $f_1, \dots, f_n$ -- набор попарно ортогональных ненулевых векторов. Тогда
    \[
    \forall x\in \mathcal{E} \hookrightarrow x=\sum_{i=1}^n \frac{(x,f_i)}{|f_i|^2}f_i
    \]
\end{proposition}
\begin{proof}
    Если набор $f^{(n)}$ -- ЛНЗ, то $\textbf{Г}=diag((f_1, f_1), \dots, (f_n, f_n))$ и $\det \textbf{Г} \neq 0$, значит $f^{(n)}$ образует базис. Следовательно, $x=\alpha_1f_1+\dots\alpha_nf_n$. Умножая данное равенство скалярно справа на $f_i$, получаем 
    \[
    \alpha_i=\frac{(x,f_i)}{|f_i|^2}
    \]
    Докажем, что набор $f^{(n)}$ всегда ЛНЗ. Рассмотрим линейную комбинацию $c_1f_1+\dots + c_n f_n=0$, умножим ее скалярно на $f_i$, получим $c_i(f_i, f_i)=0$. Но $(f_i, f_i) > 0$, значит $c_i=0$. Следовательно, $c_1=\dots=c_n=0$ и линейная комбинация является тривиальной, откуда $f^{(n)}$ -- ЛНЗ. 
\end{proof}